\chapter{Ανασκόπηση Βιβλιογραφίας και Θεωρητικό Υπόβαθρο}

\section{Εισαγωγή στην Τεχνολογία των Ψηφιακών Διδύμων (\textlatin{Digital Twins})}
Η έννοια του Ψηφιακού Διδύμου (\textlatin{Digital Twin - DT}) έχει αναδειχθεί την τελευταία δεκαετία ως μία από τις πλέον επιδραστικές τεχνολογίες στον τομέα της 4ης Βιομηχανικής Επανάστασης (\textlatin{Industry 4.0}) και του Διαδικτύου των Πραγμάτων (\textlatin{IoT}). Παρότι οι ρίζες της τεχνολογίας εντοπίζονται στον τομέα της αεροδιαστημικής (\textlatin{NASA}) για την παρακολούθηση απομακρυσμένων διαστημικών σκαφών, η ραγδαία εξέλιξη των υπολογιστικών συστημάτων και η μείωση του κόστους των αισθητήρων επέτρεψαν τη μετάβασή της στον κτιριακό και εκπαιδευτικό τομέα.

\subsection{Η Εξέλιξη: Από το Ψηφιακό Μοντέλο στο Ψηφιακό Δίδυμο}
Στη σύγχρονη ακαδημαϊκή βιβλιογραφία, αποτελεί συχνό φαινόμενο η σύγχυση και η εσφαλμένη ταύτιση της απλής τρισδιάστατης μοντελοποίησης με τα πραγματικά Ψηφιακά Δίδυμα. Προκειμένου να τεθεί ένα σαφές ερευνητικό πλαίσιο, είναι απαραίτητος ο αυστηρός εννοιολογικός διαχωρισμός των βαθμίδων ψηφιακής αναπαράστασης, ο οποίος βασίζεται αποκλειστικά στη ροή των δεδομένων (\textlatin{data flow}) μεταξύ του φυσικού και του εικονικού αντικειμένου \cite{Grieves2017}, \cite{Kritzinger2018}, \cite{Fuller2020}.

\begin{itemize}
\item \textbf{Ψηφιακό Μοντέλο (\textlatin{Digital Model}):} Αποτελεί την ψηφιακή αναπαράσταση ενός υπαρκτού ή μελλοντικού φυσικού συστήματος (όπως ένα αρχιτεκτονικό σχέδιο \textlatin{CAD} ή ένα στατικό μοντέλο \textlatin{BIM}). Η κρίσιμη παράμετρος εδώ είναι ότι δεν υφίσταται καμία απολύτως αυτοματοποιημένη ανταλλαγή δεδομένων μεταξύ του φυσικού και του ψηφιακού αντικειμένου. Κάθε αλλαγή στο ψηφιακό μοντέλο πρέπει να εισαχθεί χειροκίνητα από τον χρήστη.
\item \textbf{Ψηφιακή Σκιά (\textlatin{Digital Shadow}):} Πρόκειται για την εξέλιξη του ψηφιακού μοντέλου, όπου εγκαθιδρύεται μια αυτοματοποιημένη, μονόδρομη ροή δεδομένων από το φυσικό αντικείμενο προς το ψηφιακό. Τα περισσότερα σύγχρονα συστήματα τηλεμετρίας, όπως τα \textlatin{dashboards} που παρακολουθούν παραμέτρους (π.χ. θερμοκρασία, κατανάλωση) σε πραγματικό χρόνο, κατατάσσονται σε αυτή την κατηγορία. Μια μεταβολή στη φυσική κατάσταση (π.χ. αύξηση της θερμοκρασίας) αντικατοπτρίζεται αυτόματα στο ψηφιακό ταμπλό, αλλά το σύστημα αδυνατεί να επιδράσει αντίστροφα στο φυσικό περιβάλλον.
\item \textbf{Ψηφιακό Δίδυμο (\textlatin{Digital Twin}):} Η τρίτη και ανώτατη βαθμίδα απαιτεί την ύπαρξη μιας πλήρως αμφίδρομης, αυτοματοποιημένης ροής δεδομένων. Το ψηφιακό αντικείμενο όχι μόνο λαμβάνει ζωντανά δεδομένα από το φυσικό του αντίστοιχο (\textlatin{Sense}), αλλά είναι σε θέση να αναλύει τα δεδομένα, να λαμβάνει αποφάσεις (\textlatin{Decide}) και να στέλνει πίσω εντολές ενεργοποίησης (\textlatin{Act}), επιφέροντας αλλαγές στο φυσικό περιβάλλον (\textlatin{actuation}). Αυτός ο κλειστός βρόχος ελέγχου (\textlatin{closed-loop control}) είναι η βασική προϋπόθεση για τον χαρακτηρισμό ενός συστήματος ως αυθεντικό Ψηφιακό Δίδυμο.
\end{itemize}

Σύμφωνα με πρόσφατες συνθετικές ανασκοπήσεις της βιβλιογραφίας \cite{Han2022, Kritzinger2018}, το μέλλον των Ψηφιακών Διδύμων ξεφεύγει από την απλή οπτικοποίηση. Η ακαδημαϊκή έρευνα συγκλίνει στο ότι η επόμενη γενιά συστημάτων θα βασίζεται σε τέσσερις πυλώνες: την ενσωμάτωση φωνητικών βοηθών (\textlatin{Voice Assistants}), τη συνεχή βελτιστοποίηση υλικού και λογισμικού, την προγνωστική ανάλυση (\textlatin{Predictive Analytics}) μέσω Μηχανικής Μάθησης, με ενδεικτικές εφαρμογές όπως το \textlatin{AI Energy Forecasting}, και, τελικά, την πλήρη αυτονομία στη λήψη αποφάσεων, δηλαδή τη μετάβαση προς γνωστικά Ψηφιακά Δίδυμα (\textlatin{Cognitive Digital Twins}) \cite{Han2022}, \cite{Semeraro2021}. Στο πλαίσιο αυτό, η ανάλυση δεδομένων (\textlatin{Data Analysis}) αποκτά κεντρικό ρόλο, καθώς το \textlatin{AI} μπορεί να επεξεργάζεται πολύ μεγάλους όγκους τηλεμετρίας που συλλέγονται από το φυσικό αντικείμενο, να εντοπίζει κρυφά μοτίβα λειτουργίας και να υποστηρίζει τη βελτιστοποίηση της απόδοσής του. Για παράδειγμα, εάν το ψηφιακό δίδυμο ενός βιομηχανικού ψυγείου εντοπίσει ασυνήθιστους κραδασμούς ή θερμικές διακυμάνσεις, οι αλγόριθμοι \textlatin{AI} μπορούν να αναλύσουν τα σχετικά δεδομένα και να προβλέψουν πιθανές βλάβες πριν αυτές συμβούν, ενισχύοντας πρακτικά την προγνωστική συντήρηση. Παράλληλα, η διαδραστική επικοινωνία (\textlatin{Interactive Communication}) μετασχηματίζει τη σχέση του χρήστη με την πλατφόρμα, αφού αντί για αποκλειστική ανάγνωση στατικών γραφημάτων καθίσταται εφικτή η αλληλεπίδραση μέσω καθημερινής γλώσσας με χρήση τεχνικών επεξεργασίας φυσικής γλώσσας (\textlatin{Natural Language Processing - NLP}). Η προοπτική αυτή έχει ιδιαίτερη σημασία για τις σχολικές και κτιριακές εφαρμογές, καθώς μετατοπίζει το κέντρο βάρους από την παθητική αναπαράσταση του χώρου προς την ενεργή υποστήριξη της λειτουργίας, της βελτιστοποίησης και της επιχειρησιακής διαχείρισής του.

\subsection{Η Αρχιτεκτονική Πέντε Επιπέδων (\textlatin{5-Layer Architecture})}
Για την επιτυχή υλοποίηση ενός Ψηφιακού Διδύμου σε πολύπλοκα περιβάλλοντα, όπως είναι ένα σχολικό συγκρότημα (\textlatin{Smart Campus}), Η αρχιτεκτονική πέντε επιπέδων που υιοθετείται στην παρούσα εργασία βασίζεται στη σύνθεση των προτύπων που προτείνουν οι Han et al. \cite{Han2022} για έξυπνες πανεπιστημιουπόλεις και οι Fuller et al. \cite{Fuller2020} για Ψηφιακά Δίδυμα γενικότερα, αξιοποιώντας την ιεράρχηση \textlatin{Physical, Connection, Data, Virtual} και \textlatin{Service layers} με αμφίδρομη ροή \textlatin{Sense~$\uparrow$ / Act~$\downarrow$}. Το μοντέλο αυτό διασφαλίζει την επεκτασιμότητα (\textlatin{scalability}) και τη διαλειτουργικότητα (\textlatin{interoperability}) των συστημάτων. Τα επίπεδα αυτά περιλαμβάνουν: (1) το \textit{Φυσικό Επίπεδο (\textlatin{Physical Layer})} όπου συλλέγονται τα πρωτογενή δεδομένα μέσω αισθητήρων IoT, (2) το \textit{Επίπεδο Δικτύου (\textlatin{Connection Layer})} που εξασφαλίζει την ασφαλή διακίνηση δεδομένων (π.χ. μέσω \textlatin{MQTT}), (3) την \textit{Πλατφόρμα Δεδομένων (\textlatin{Data Layer})} για την ενοποίηση και αποθήκευση σε βάσεις χρονοσειρών, (4) το \textit{Εικονικό Επίπεδο (\textlatin{Virtual Layer})} που παρέχει την τρισδιάστατη χωρική οπτικοποίηση (\textlatin{3D Models}), και τέλος, (5) το \textit{Επίπεδο Υπηρεσιών (\textlatin{Service Layer})} όπου εκτελούνται οι αλγόριθμοι λήψης αποφάσεων (\textlatin{Rules Engine}) και οι αυτοματισμοί. Η συγκεκριμένη αρχιτεκτονική υιοθετήθηκε πλήρως ως ο κατευθυντήριος άξονας για την ανάπτυξη της πλατφόρμας της παρούσας διπλωματικής εργασίας.

\section{Έξυπνες Πανεπιστημιουπόλεις (\textlatin{Smart Campuses}) και το Ευρωπαϊκό Πλαίσιο}
Η εφαρμογή των Ψηφιακών Διδύμων έχει αρχίσει να μετατοπίζεται από την αμιγώς βιομηχανική παραγωγή (\textlatin{Industry 4.0}) στον αστικό και κτιριακό τομέα, στοχεύοντας στη δημιουργία Έξυπνων Πόλεων (\textlatin{Smart Cities}) και Έξυπνων Σχολικών Υποδομών (\textlatin{Smart Campuses}).

\subsection{Ευρωπαϊκές Πολιτικές και Πρωτοβουλίες}
Η υιοθέτηση των ψηφιακών διδύμων ευθυγραμμίζεται σε μεγάλο βαθμό με τους στρατηγικούς στόχους της Ευρωπαϊκής Ένωσης. Το πρόγραμμα \textit{DIGITAL EUROPE} δίνει ρητή έμφαση στη δημιουργία ``Τοπικών Ψηφιακών Διδύμων'' (Local Digital Twins) ως βασικό εργαλείο για την ενίσχυση της βιωσιμότητας των έξυπνων κοινοτήτων και την αποτελεσματικότερη διαχείριση κτιριακών υποδομών \cite{EuropeanCommission2023}. Επιπρόσθετα, η Ευρωπαϊκή Πράσινη Συμφωνία καθιστά επιτακτική την ενεργειακή αναβάθμιση του κτιριακού αποθέματος \cite{EUParliament2024}. Η ενσωμάτωση τεχνολογιών \textlatin{IoT} σε σχολεία αποτελεί ένα καίριο βήμα για τη μετάβαση σε υποδομές μηδενικών εκπομπών.
Η υιοθέτηση των ψηφιακών διδύμων ευθυγραμμίζεται σε μεγάλο βαθμό με τους στρατηγικούς στόχους της Ευρωπαϊκής Ένωσης. Το πρόγραμμα \textit{DIGITAL EUROPE} δίνει ρητή έμφαση στη δημιουργία ``Τοπικών Ψηφιακών Διδύμων'' (\textlatin{Local Digital Twins}) ως βασικό εργαλείο για την ενίσχυση της βιωσιμότητας των έξυπνων κοινοτήτων και την αποτελεσματικότερη διαχείριση κτιριακών υποδομών \cite{EuropeanCommission2023}. Επιπρόσθετα, η Ευρωπαϊκή Πράσινη Συμφωνία καθιστά επιτακτική την ενεργειακή αναβάθμιση του κτιριακού αποθέματος \cite{EUParliament2024}. Η ενσωμάτωση τεχνολογιών \textlatin{IoT} σε σχολεία αποτελεί ένα καίριο βήμα για τη μετάβαση σε υποδομές μηδενικών εκπομπών.
Η υιοθέτηση των ψηφιακών διδύμων ευθυγραμμίζεται σε μεγάλο βαθμό με τους στρατηγικούς στόχους της Ευρωπαϊκής Ένωσης. Το πρόγραμμα \textit{\textlatin{DIGITAL EUROPE}} δίνει ρητή έμφαση στη δημιουργία ``Τοπικών Ψηφιακών Διδύμων'' (\textlatin{Local Digital Twins}) ως βασικό εργαλείο για την ενίσχυση της βιωσιμότητας των έξυπνων κοινοτήτων και την αποτελεσματικότερη διαχείριση κτιριακών υποδομών \cite{EuropeanCommission2023}. Επιπρόσθετα, η Ευρωπαϊκή Πράσινη Συμφωνία καθιστά επιτακτική την ενεργειακή αναβάθμιση του κτιριακού αποθέματος \cite{EUParliament2024}. Η ενσωμάτωση τεχνολογιών \textlatin{IoT} σε σχολεία αποτελεί ένα καίριο βήμα για τη μετάβαση σε υποδομές μηδενικών εκπομπών.
Η υιοθέτηση των ψηφιακών διδύμων ευθυγραμμίζεται σε μεγάλο βαθμό με τους στρατηγικούς στόχους της Ευρωπαϊκής Ένωσης. Το πρόγραμμα \textit{\textlatin{DIGITAL EUROPE}} δίνει ρητή έμφαση στη δημιουργία ``Τοπικών Ψηφιακών Διδύμων'' (\textlatin{Local Digital Twins}) ως βασικό εργαλείο για την ενίσχυση της βιωσιμότητας των έξυπνων κοινοτήτων και την αποτελεσματικότερη διαχείριση κτιριακών υποδομών \cite{EuropeanCommission2023}. Επιπρόσθετα, η Ευρωπαϊκή Πράσινη Συμφωνία καθιστά επιτακτική την ενεργειακή αναβάθμιση του κτιριακού αποθέματος \cite{EUParliament2024}. Η ενσωμάτωση τεχνολογιών \textlatin{IoT} σε σχολεία αποτελεί ένα καίριο βήμα για τη μετάβαση σε υποδομές μηδενικών εκπομπών.

\subsection{Συναφές Ερευνητικό Έργο και Πιλοτικές Εφαρμογές}
Σε ευρωπαϊκό και εθνικό επίπεδο, αρκετές πρόσφατες μελέτες έχουν εξετάσει τις προοπτικές των ψηφιακών διδύμων. Μια χαρακτηριστική ελληνική περίπτωση είναι η μελέτη των \textlatin{Evangelou et al.} \cite{Evangelou2022}, οι οποίοι διερεύνησαν την ενσωμάτωση Μοντέλων Πληροφοριών Κτιρίου (\textlatin{BIM}) με Γεωγραφικά Συστήματα Πληροφοριών (\textlatin{GIS}) για τη δημιουργία Ψηφιακών Διδύμων σε αστικό επίπεδο στην Ελλάδα, αναδείκνύοντας τη σημασία της χωρικής τεκμηρίωσης στη διαχείριση έξυπνων υποδομών και τη λήψη αποφάσεων σε επίπεδο διαχείρισης εγκαταστάσεων (\textlatin{Facility Management}).

Σε διεθνές επίπεδο, η έρευνα των \textlatin{Cui et al.} (2023) επικεντρώθηκε στην "Ψηφιακή προσομοίωση διδύμων ενός σχολικού κτιρίου στη Δανία", αποδεικνύοντας την αξία προηγμένων αλγορίθμων \textlatin{Model Predictive Control} (\textlatin{MPC}) για τη βελτιστοποίηση ενεργειακής χρήσης σε σχολικά κτίρια \cite{Cui2023}, προσέγγιση διαφορετική από την χαμηλού κόστους \textlatin{COTS} αρχιτεκτονική της παρούσας εργασίας. Η εν λόγω μελέτη απέδειξε την αξία της χρήσης προηγμένων αλγορίθμων \textlatin{Model Predictive Control} (\textlatin{MPC}) και γραμμικών μοντέλων χρονοσειρών για την προσομοίωση και βελτιστοποίηση της ενεργειακής συμπεριφοράς σχολικών κτιρίων, λαμβάνοντας υπόψη εξωτερικές μεταβλητές όπως η ηλιακή ακτινοβολία και η θερμική αδράνεια της κατασκευής \cite{Cui2023}. Παράλληλα, πρόσφατες υλοποιήσεις σε σχολικά περιβάλλοντα, όπως το ερευνητικό έργο «\textlatin{SchoolAIR}» στην Πορτογαλία, ανέδειξαν τη βιωσιμότητα συστημάτων χαμηλού κόστους (βασισμένων σε μικροελεγκτές \textlatin{ESP32} και \textlatin{Raspberry Pi}) για τη διαχείριση της ποιότητας του εσωτερικού αέρα (\textlatin{Indoor Air Quality --- IAQ}) \cite{Barros2024}.

Αν και τέτοιου είδους πρωτοβουλίες επιβεβαιώνουν ότι η παρακολούθηση (\textlatin{monitoring}) με εμπορικά διαθέσιμο (\textlatin{COTS - Commercial Off-The-Shelf}) εξοπλισμό είναι απολύτως εφικτή, συνήθως περιορίζονται στη βαθμίδα της «Ψηφιακής Σκιάς» (\textlatin{Digital Shadow}). Παρέχουν, δηλαδή, δεδομένα σε πραγματικό χρόνο αλλά δεν κλείνουν τον βρόχο ελέγχου με αυτοματοποιημένες ενέργειες στον φυσικό χώρο. Σε ακαδημαϊκό επίπεδο, οι \textlatin{Balla et al.} \cite{Balla2023} τεκμηρίωσαν την αξία της αμφίδρομης ροής δεδομένων σε εκπαιδευτικά σενάρια ψηφιακών διδύμων, αποδεικνύοντας ότι η πλήρης υλοποίηση του κύκλου \textlatin{Sense--Decide--Act} απαιτεί σαφή αρχιτεκτονικό διαχωρισμό μεταξύ παθητικής παρακολούθησης (\textlatin{Digital Shadow}) και ενεργού παρέμβασης (\textlatin{Digital Twin}). Αντίστοιχα, ευρύτερες προσεγγίσεις που αξιοποιούν συστήματα \textlatin{BIM-GIS} σε μεγάλες πανεπιστημιουπόλεις για την κεντρική διαχείριση παγίων (\textlatin{Asset Management}) αποδεικνύουν τη σπουδαιότητα της τρισδιάστατης οπτικοποίησης, υπογραμμίζοντας ταυτόχρονα την τεχνική πολυπλοκότητα και το υψηλό κόστος μετάβασης σε πλήρως αμφίδρομο Ψηφιακό Δίδυμο εντός υφιστάμενων υποδομών \cite{Evangelou2022}, \cite{Fuller2020}.

Στον ευρύτερο τομέα των έξυπνων κτιριακών υποδομών, ερευνητές έχουν εξερευνήσει τη σύνδεση των Ψηφιακών Διδύμων με τεχνολογίες Εκτεταμένης Πραγματικότητας (\textlatin{Extended Reality --- XR}). Οι \textlatin{Li et al.} \cite{Li2024} απέδειξαν ότι η ενσωμάτωση \textlatin{XR workflows} σε περιβάλλοντα \textlatin{smart building operations} επιτρέπει τη δημιουργία Εικονικών Εργαστηρίων (\textlatin{Virtual Labs}) για εκπαιδευτικούς και τεχνικούς χρήστες, ανοίγοντας τον δρόμο για πιο διαδραστικές μορφές επιχειρησιακής διαχείρισης.

\section{Ποιότητα Εσωτερικού Περιβάλλοντος (\textlatin{IEQ}) και Μαθησιακή Διαδικασία}
Ένας από τους κύριους πυλώνες στους οποίους εστιάζει ένα σχολικό Ψηφιακό Δίδυμο είναι η βελτιστοποίηση της Ποιότητας Εσωτερικού Περιβάλλοντος (\textlatin{Indoor Environmental Quality - IEQ}). Το σχολικό περιβάλλον διαδραματίζει καθοριστικό ρόλο στη συγκέντρωση, τη γνωστική απόδοση και την υγεία των μαθητών και των εκπαιδευτικών.

\subsection{Η Ακουστική Άνεση ως Κρίσιμος Παράγοντας}

Σύμφωνα με τις επίσημες κατευθυντήριες οδηγίες του Παγκόσμιου Οργανισμού Υγείας (\textlatin{WHO}), η μέση στάθμη θορύβου υποβάθρου (\textlatin{background noise}) σε μια άδεια (\textlatin{unoccupied}) σχολική αίθουσα δεν θα πρέπει να υπερβαίνει τα 35~\textlatin{dBA} \cite{WHO1999}, \cite{Kephalopoulos2015}. Ο στόχος αυτού του αυστηρού ορίου είναι η διασφάλιση επαρκούς Λόγου Σήματος προς Θόρυβο (\textlatin{Signal-to-Noise Ratio - SNR}) της τάξης των +15 \textlatin{dB}, ώστε η φωνή του εκπαιδευτικού να είναι απολύτως καθαρή και κατανοητή από κάθε σημείο της αίθουσας. Παράλληλα, οργανισμοί όπως ο \textlatin{ANSI} \cite{ANSI2010} και ο \textlatin{ISO} \cite{ISO2008} ορίζουν αντίστοιχα όρια για ακουστικές παραμέτρους σχολικών χώρων, συμπεριλαμβανομένης της μέγιστης αντήχησης (\textlatin{RT60}), θέτοντας το ευρύτερο κανονιστικό πλαίσιο εντός του οποίου εντάσσεται η παρακολούθηση της ισοδύναμης συνεχούς στάθμης ηχητικής πίεσης (\textlatin{LAeq}) της παρούσας μελέτης.

Ωστόσο, κατά τη διάρκεια ενός ενεργού μαθήματος, τα επίπεδα θορύβου είναι εξαιρετικά δυναμικά. Η παρούσα βιβλιογραφία υποδεικνύει ότι απαιτούνται δυναμικοί μηχανισμοί ειδοποίησης, χωρίς όμως αυτοί να γίνονται παρεμβατικοί. Η χρήση ισχυρών ηχητικών συναγερμών (\textlatin{buzzers}) σε περιπτώσεις φασαρίας έχει αποδειχθεί λανθασμένη πρακτική, καθώς προσθέτει επιπλέον θόρυβο στο ήδη βεβαρημένο περιβάλλον, διακόπτοντας το μάθημα. Αντιθέτως, η χρήση οπτικών ενδείξεων (όπως η "\textlatin{quiet-lamp}" που αναπτύχθηκε σε αυτή τη διπλωματική) αποτελεί μια προσέγγιση παθητικής αυτορρύθμισης συμπεριφοράς (\textlatin{behavioral nudging}) για τη διατήρηση του θορύβου κάτω από δυναμικά κατώφλια λειτουργίας \cite{Dockrell2006}. Η επιλογή οπτικής αντί ηχητικής ανατροφοδότησης ελαχιστοποιεί την πρόσθετη ακουστική επιβάρυνση της αίθουσας \cite{Dockrell2006}, σύμφωνα με τα ευρήματα της βιβλιογραφίας για την επίδραση του θορύβου στη μαθησιακή διαδικασία \cite{MogasRecalde2021}.

\subsection{Θερμική Άνεση και Ποιότητα Αέρα}
Πέραν της ακουστικής, η θερμική άνεση και η συγκέντρωση Διοξειδίου του Άνθρακα (\textlatin{CO2}) είναι κρίσιμες μεταβλητές. Σύμφωνα με μελέτες που αφορούν το μικροκλίμα των σχολείων, συγκεντρώσεις \textlatin{CO2} άνω των 700--1000~\textlatin{ppm} έχουν συσχετιστεί με μειωμένη γνωστική απόδοση και αίσθημα υπνηλίας σε εσωτερικούς χώρους \cite{Allen2016}. Η ένταξη αισθητήρων που παρακολουθούν συνεχώς τη θερμοκρασία, την υγρασία και τον αερισμό στο Ψηφιακό Δίδυμο, επιτρέπει στη διοίκηση να εξασφαλίζει τις ιδανικές συνθήκες για μάθηση, προλαμβάνοντας καταστάσεις δυσφορίας πριν αυτές επηρεάσουν τους μαθητές.

\section{Ενεργειακή και Λειτουργική Διαχείριση (\textlatin{Facility Management})}
Ο δεύτερος κεντρικός πυλώνας της βιβλιογραφίας των ψηφιακών διδύμων αφορά τη Λειτουργική Διαχείριση των Εγκαταστάσεων (\textlatin{Facility Management}) και την Ενεργειακή Αποδοτικότητα. Τα σχολικά συγκροτήματα χαρακτηρίζονται από ένα ιδιαίτερο προφίλ λειτουργίας: παρουσιάζουν υψηλή ζήτηση ενέργειας κατά τις πρωινές ώρες, ενώ παραμένουν ανενεργά για μεγάλα διαστήματα (απογεύματα, Σαββατοκύριακα, διακοπές).

\subsection{Η Αντιμετώπιση των Κρυφών Φορτίων (\textlatin{Phantom Loads})}
Παρά την απουσία δραστηριότητας κατά τις ώρες αδράνειας, η βιβλιογραφία καταδεικνύει ότι τα δημόσια και ιδιωτικά κτίρια σπαταλούν σημαντικά ποσά ενέργειας λόγω των λεγόμενων "\textlatin{κρυφών φορτίων}" (\textlatin{phantom loads}). Εξοπλισμός σε κατάσταση αναμονής (\textlatin{stand-by}), φώτα που ξεχάστηκαν ανοιχτά, και κυρίως βαρέα συστήματα ψύξης-θέρμανσης (όπως οι επαγγελματικοί καταψύκτες και τα ψυγεία των σχολικών κυλικείων) λειτουργούν άσκοπα.

Η μετάβαση στο Ψηφιακό Δίδυμο προσφέρει μια άμεση, τεχνολογική λύση σε αυτό το πρόβλημα. Με την εγκατάσταση έξυπνων μετρητών κατανάλωσης (όπως τα ρελέ της \textlatin{Shelly}) σε συνδυασμό με αισθητήρες κατάστασης (π.χ. μαγνητικές επαφές θυρών), το σύστημα μπορεί να ανιχνεύσει λειτουργικές ανωμαλίες. Η εφαρμογή κανόνων (\textlatin{rules engine}) που ειδοποιούν άμεσα όταν μια πόρτα ψυγείου παραμένει ανοιχτή αποτρέπει την απώλεια ψύξης και ελαχιστοποιεί την άσκοπη λειτουργία των συμπιεστών \cite{Apostolopoulos2022}. Η παρακολούθηση ενεργειακών φορτίων σε πραγματικό χρόνο αποτελεί βασικό εργαλείο για τη βελτιστοποίηση της λειτουργίας σχολικών κτιρίων \cite{Fokaides2020}, \cite{Markoska2019}.

\subsection{Προγνωστική Συντήρηση (\textlatin{Predictive Maintenance})}
Ένα από τα κορυφαία οφέλη που αναγνωρίζεται στη βιβλιογραφία των \textlatin{Digital Twins} είναι η δυνατότητα Προγνωστικής Συντήρησης. Μέσω της συνεχούς καταγραφής των ενεργειακών προφίλ (π.χ. της διαρκούς καταγραφής της Ισχύος σε \textlatin{Watts}), το σύστημα μπορεί να αναγνωρίσει εγκαίρως μικρές αποκλίσεις που υποδηλώνουν μηχανική φθορά. Εάν ο συμπιεστής ψυγείου εμφανίσει παρατεταμένη αύξηση κατανάλωσης ισχύος πέρα από τα αναμενόμενα επίπεδα λειτουργίας, το σύστημα εντοπίζει την ανωμαλία και ειδοποιεί τους τεχνικούς για ενδεχόμενη μηχανική βλάβη ή μειωμένη απόδοση του ψυκτικού εξοπλισμού \cite{Cabrera2020}, \cite{Han2022}. Η προσέγγιση αυτή αποτελεί πρακτική εφαρμογή της Προγνωστικής Συντήρησης (\textlatin{Predictive Maintenance}) σε επίπεδο σχολικής εγκατάστασης.

\section{Το Ερευνητικό Κενό και η Συνεισφορά της Πλατφόρμας}
Αναλύοντας το υφιστάμενο ακαδημαϊκό και τεχνολογικό τοπίο, διαπιστώνεται ότι, ενώ οι θεωρητικές προσεγγίσεις για τα \textlatin{Smart Campuses} αφθονούν, οι πρακτικές, πλήρως λειτουργικές υλοποιήσεις \textit{αληθινών} Ψηφιακών Διδύμων στην Πρωτοβάθμια και Δευτεροβάθμια εκπαίδευση παραμένουν σπάνιες. Στη σύγχρονη βιβλιογραφία, η πλειονότητα των σχολικών εγκαταστάσεων \textlatin{IoT} ταξινομείται ως \textlatin{Digital Shadow} \cite{Kritzinger2018}: λειτουργούν δηλαδή με μονόδρομη ροή δεδομένων, προσφέροντας εξαιρετικά δεδομένα παρακολούθησης και ειδοποιήσεις μέσω \textlatin{dashboards}, στερούνται όμως της ικανότητας αμφίδρομου ελέγχου (\textlatin{closed-loop control}) που απαιτείται για τον χαρακτηρισμό τους ως πλήρες \textlatin{Digital Twin} \cite{Grieves2017}, \cite{Fuller2020}.

Παραδοσιακά, η επίτευξη αυτού του επιπέδου ελέγχου ανατίθεται σε βιομηχανικά Συστήματα Διαχείρισης Κτιρίων (\textlatin{BMS}). Ωστόσο, τα συστήματα αυτά απαιτούν σημαντικά υψηλότερο αρχικό κόστος κεφαλαίου (\textlatin{CAPEX}) σε σχέση με λύσεις ανοιχτού κώδικα, βασίζονται σε κλειστά και ιδιόκτητα πρωτόκολλα (\textlatin{proprietary hardware}) και προϋποθέτουν εξειδικευμένο προσωπικό συντήρησης, καθιστώντας τα πρακτικά απρόσιτα για το μέσο σχολικό συγκρότημα. Στον αντίποδα, οι φθηνές εμπορικές λύσεις έξυπνου σπιτιού (\textlatin{smart home apps}) λειτουργούν ως μεμονωμένες «νησίδες» αυτοματισμού, παρέχοντας περιορισμένες δυνατότητες χωρίς ενοποιημένη διαχείριση επιπέδου εγκατάστασης (\textlatin{Enterprise Level}).

Το ερευνητικό κενό που έρχεται να καλύψει η παρούσα διπλωματική εργασία εντοπίζεται ακριβώς στη γεφύρωση αυτών των δύο άκρων: στην ανάπτυξη μιας \textbf{επαναλήψιμης (\textlatin{replicable}), χαμηλού κόστους μεθοδολογίας, αυστηρά βασισμένης σε τεχνολογίες ανοιχτού κώδικα (\textlatin{open-source})}. Συνδυάζοντας προσιτούς μικροελεγκτές (όπως το \textlatin{ESP32}) και συσκευές \textlatin{Edge} (\textlatin{Shelly, Sonoff}) με βιομηχανικά πρωτόκολλα μηνυμάτων (\textlatin{MQTT}) και ένα ισχυρό σημασιολογικό ενδιάμεσο λογισμικό (\textlatin{semantic middleware - openHAB}), η αναπτυχθείσα πλατφόρμα εγκαθιδρύει την πλήρη αλληλουχία \textit{\textlatin{Sense - Decide - Act}}. Αποδεικνύει έτσι στην πράξη ότι ο αμφίδρομος έλεγχος, η τρισδιάστατη οπτικοποίηση και η μετάβαση σε ένα ολοκληρωμένο Έξυπνο Σχολείο είναι τεχνολογικά εφικτή και οικονομικά βιώσιμη, μετασχηματίζοντας τον σχολικό χώρο σε έναν ενεργό και δυναμικό ψηφιακό οργανισμό. Επιπλέον, η πλατφόρμα σχεδιάστηκε με ενσωματωμένη βιβλιοθήκη προτύπων χώρων (\textlatin{template library}), η οποία επιτρέπει την επέκτασή της σε νέες εκπαιδευτικές εγκαταστάσεις με ελάχιστη προσαρμογή, ενισχύοντας τον επαναλήψιμο (\textlatin{replicable}) χαρακτήρα της λύσης.

\begin{table}[htbp]
\centering
\caption{Συγκριτική ανάλυση παραδοσιακής διαχείρισης έναντι Ψηφιακού Διδύμου}
\label{tab:traditional_vs_dt}
\begin{tabular}{|l|l|l|l|}
\hline
\textbf{Μετρική} & \textbf{Παραδοσιακή Διαχείριση} & \textbf{Ψηφιακό Δίδυμο} & \textbf{Βελτίωση} \\
\hline
Συχνότητα Παρακολούθησης & 1 μέτρηση/ώρα (χειροκίνητη) & 100+ μετρήσεις/ώρα (αυτόματη) & $\times$100 ταχύτερα \\
\hline
Χρόνος Ανίχνευσης Βλάβης & $\sim$48 ώρες & $\sim$30 λεπτά & $\sim$96$\times$ ταχύτερα \\
\hline
Ακρίβεια Δεδομένων & 70–80\% (ανθρώπινο σφάλμα) & 98,5\% (επαληθευμένη) & +31\% \\
\hline
Κατανάλωση Ενέργειας & 223 \textlatin{kWh/ημέρα} & 179 \textlatin{kWh/ημέρα} & -19,7\% \\
\hline
Διαθεσιμότητα Εποπτείας & Κατά τις ώρες εργασίας & 99,95\% \textlatin{uptime} (24/7) & Συνεχής \\
\hline
Κόστος Ανίχνευσης Αστοχίας & Υψηλό (μετά τη βλάβη) & Χαμηλό (προληπτικό) & \textlatin{Predictive} \\
\hline
\end{tabular}
\end{table}

\noindent\textit{Σημείωση: Τα δεδομένα του πίνακα προκύπτουν από τη σύγκριση της κατάστασης πριν (Σεπ. 2024 -- Φεβ. 2025) και μετά (Μάρ. 2025 -- Νοε. 2025) την εγκατάσταση της πλατφόρμας στα Εκπαιδευτήρια Πλάτων.}
